\documentclass[french]{book}
\usepackage{teach}
\usepackage{verbatim}
\usepackage{amsmath,amsfonts,amssymb}
\usepackage{pgfplots}
\usepackage{tikz}
\usepackage{mwe}
\RequirePackage{graphicx}
\RequirePackage{ifpdf}
 \ifpdf
   \DeclareGraphicsRule{*}{mps}{*}{}
 \fi




\RequirePackage{fancyhdr}
\RequirePackage{titlesec}

  \usepackage{fancyhdr}

  \usepackage{amsmath}
  \usepackage{amssymb}
  \usepackage{amsfonts}
  \usepackage{amsthm}
  \usepackage{mathrsfs}

  \usepackage{array}
  \usepackage{multicol}
  \setlength{\multicolsep}{3pt}

  \usepackage{graphicx}
  \graphicspath{{Figures/}}
  \usepackage{pstricks,pst-plot,pst-text,pst-tree,pst-eucl}
  \usepackage[all]{xy}

  \usepackage{fancybox}
  \usepackage{color}
    \usepackage{stmaryrd}
\usepackage{etex}
%\usepackage{preambuleTrm}
%\usepackage{enumitem}
%\usepackage{tikz}
%
\usetikzlibrary{calc}
%\usepackage[xcas]{pro-courbes}
%\usepackage[autolanguage]{numprint}
\usepackage[french,vlined]{algorithm2e}
\usepackage{cclicenses}
%\usepackage{cclicence}
%\usepackage{docacompleter}

\newcommand{\cacheb}[1]{#1}
\newcommand{\cacheg}[1]{#1} 
\newcommand{\cacher}[1]{#1}
\newcommand{\cachegg}[1]{#1} 

% \newcommand{\cacheb}[1]{\dotfill\textcolor{white}{#1}\textcolor{white}{#1}}
% \newcommand{\cacheg}[1]{\dotfill\textcolor{0.9white}{#1}\textcolor{0.9white}{#1}} 

% \newcommand{\cachegg}[1]{\dots\dots\textcolor{0.8white}{#1}\textcolor{0.8white}{#1}} 
\newcolumntype{A}{>{\arraybackslash}X}

\usepackage{textcomp}
\usepackage{mathtools}
%\usepackage[xcas]{pro-statdiv}
%  \usepackage[frenchb]{babel}
  \usepackage{hyperref}

\usepackage{subfig}
\pgfplotsset{compat=newest}
\usepackage[all]{xy}
%\usepackage{geometry}
%\geometry{top=1cm, bottom=1cm, left=2cm, right=2cm}
\usepackage[utf8]{inputenc}
\usepackage[french]{babel}
\redefinecolor{thm}{title}{bgcolor}{alizarin}
\redefinecolor{prop}{title}{bgcolor}{meatbrown}
\redefinecolor{defn}{title}{bgcolor}{olivine}
\definecolor{alizarin}{rgb}{0.82, 0.1, 0.26}
\definecolor{meatbrown}{rgb}{0.9, 0.72, 0.23}
\definecolor{olivine}{rgb}{0.6, 0.73, 0.45}
\usepackage{mathrsfs}
%%
\newcommand{\R}{\mathbb {R}}
%\newcommand{\C}{\mathds {C}}
\newcommand{\N}{\mathbb {N}}
\newcommand{\Z}{\mathbb {Z}}
\newcommand{\Q}{\mathbb {Q}}
\newcommand{\D}{\mathbb {D}}
%%
\newcommand{\se}{\geq}
\newcommand{\ie}{\leq}
\newcommand{\tm}{\times}
\newcommand\binomialCoefficient[2]{%
    % Store values 
    \c@pgf@counta=#1% n
    \c@pgf@countb=#2% k
    %
    % Take advantage of symmetry if k > n - k
    \c@pgf@countc=\c@pgf@counta%
    \advance\c@pgf@countc by-\c@pgf@countb%
    \ifnum\c@pgf@countb>\c@pgf@countc%
        \c@pgf@countb=\c@pgf@countc%
    \fi%
    %
    % Recursively compute the coefficients
    \c@pgf@countc=1% will hold the result
    \c@pgf@countd=0% counter
    \pgfmathloop% c -> c*(n-i)/(i+1) for i=0,...,k-1
        \ifnum\c@pgf@countd<\c@pgf@countb%
        \multiply\c@pgf@countc by\c@pgf@counta%
        \advance\c@pgf@counta by-1%
        \advance\c@pgf@countd by1%
        \divide\c@pgf@countc by\c@pgf@countd%
    \repeatpgfmathloop%
    \the\c@pgf@countc%
}
\makeatother
%\newcommand{\ell}{l}
%\newcommand{\ell'}{l'}
\newcommand{\ds}{\displaystyle}
\newcommand{\limn}{\lim_{n\rightarrow +\infty}}
\newcommand{\limo}[1][x]{\ds\lim_{#1\rightarrow 0}}
\newcommand{\limite}[2][x]{\ds\lim_{#1\rightarrow #2}}
\newcommand{\limpinf}[1][x]{\ds\lim_{#1 \rightarrow +\infty}}
\newcommand{\limminf}[1][x]{\ds\lim_{#1\rightarrow -\infty}}
\newcommand{\limiteg}[2][x]{\ds\lim_{#1\xrightarrow{<}#2}}
\newcommand{\limited}[2][x]{\ds\lim_{#1\xrightarrow{>}#2}}
\newcommand{\anglevec}[2]{(\overrightarrow{#1};\overrightarrow{#2})}
\newcommand{\co}[2]{C_#2^#1}
\newcommand{\ssi}{\Longleftrightarrow}
%%
\newcommand{\vect}[1]{\overrightarrow{#1}}
\newcommand{\oij}{(\textit{O},{\overrightarrow{i}},{\overrightarrow{j}})}
%
\newcommand{\cp}[2]{%
        \begin{pmatrix}
        #1\\
        #2
        \end{pmatrix}%
        }
%
\newcommand{\calb}{\mathscr{B}}
\newcommand{\calc}{\mathscr{C}}
\newcommand{\cald}{\mathscr{D}}
\newcommand{\cale}{\mathscr{E}}
\newcommand{\calf}{\mathscr{F}}
\newcommand{\calp}{\mathscr{P}}
\newcommand{\fr}[2]{\dfrac{#1}{#2}}
\newcommand{\cals}{\mathscr{S}}
\newcommand{\calt}{\mathscr{T}}
\newcommand{\ve}[1]{\overrightarrow{#1}}
\newcommand{\pa}[1]{(#1)}
\newcommand{\ps}[2]{\overrightarrow{#1}.\overrightarrow{#2}}
\newcommand{\bbr}{\mathbb{R}}
\newcommand{\al}{\alpha}
\newcommand{\la}{\lambda}
\newcommand{\DR}{\cald}
\newcommand{\PR}{\calp}
\newcommand{\nor}[1]{ \Vert #1 \Vert}
\newcommand{\abv}[1]{\vert #1 \vert}
\newcommand{\coordpp}[3]{\begin{pmatrix}
#1\\
#2\\
#3
\end{pmatrix}}
%
\newcommand{\suite}[1][u]{\left( #1_{n} \right)_{n\in\N}}
\newcommand{\suitep}[1][u]{\left( #1_{n} \right)_{n\in\N^{*}}}
\usetikzlibrary{arrows}


\begin{document}



\chapter{Logarithme népérien}


\section{Des fonctions vérifiant $f(a \times b)=f(a)+
f(b)$}

 Historiquement, le logarithme népérien a été pour la
première fois mis en évidence par l'\'Ecossais John
\textsc{Napier} (en français Jean Néper..) au tout début du
$XVII^{e}$ siècle. Afin de faciliter la vie des astronomes,
navigateurs, financiers de l'époque qui étaient confrontés à des
calculs...astronomiques, John rechercha une fonction qui puisse
transformer des produits très compliqués à calculer en sommes plus
abordables. Il a donc été amené à résoudre une \textit{équation
fonctionnelle}, $i.e.$ il a recherché les fonctions $f$ vérifiant
$f(a\times b)=f(a)+f(b)$

\begin{comment}
Il a pu alors établir des Tables de logarithmes, complétées au fil
des ans par des potes mathématiciens (nous verrons comment ils ont
pu se débrouiller en exercice). À partir de deux nombres $a$ et
$b$, on lit sur les Tables leurs logarithmes $\ln a$ et $\ln b$;
on calcule facilement $\ln a+ \ln b$ qui est égal à $\ln ab$, puis
on cherche sur les Tables le nombre qui admet pour logarithme $\ln
ab$ et qui est bien sûr $ab$.

\end{comment}
\begin{comment}
\subsection{Existe-t-il des primitives de x $\mapsto$ 1/x~?}

Autre problème~: nous connaissons des primitives des fonctions qui
à $x$ associent respectivement $x^2$, $x^1$, $x^0$, $x^{-2}$,
$x^{-3}$, etc. Vous avez tout de suite remarqué que nous avons
oublié quelqu'un~: quelles peuvent être les primitives de la
fonction qui à $x$ associe $x^{-1}=1/x$~? Une rapide enquête
mathématique nous conduit à trouver qu'il s'agit en fait de
fonctions vérifiant l'équation fonctionnelle du copain John. Nous
le verifierons là encore en exercice, et c'était l'approche au
programme jusqu'à l'année dernière.

\subsection{Quelle est la réciproque de la fonction
exponentielle~?}


Cette  année, comme  d'habitude, nous  roulons pour  le  nucléaire. Nous
avons  déjà   défini  une  fonction   obéissant  à  la   \textsl{loi  de
décomposition radioactive}, à savoir  la fonction exponentielle. Mais de
nouveaux  problèmes  se posent  au  moment  de  construire de  nouvelles
centrales~:  dans  combien  de   milliards  d'années  les  habitants  de
Tchernobyl ne  risqueront plus d'attraper  un cancer de la  thyroïde en
ingérant les légumes irradiés de  leurs potagers. Le Physicien est alors
amené à résoudre  une équation d'inconnue $y$ du  type $e^y=32$ Quel est
donc  ce  $y$  dont  l'exponentielle vaut  32~?  Fabrice  \textsc{Couenne},
célèbre physicien du  $XXI^e$ siècle, vous a déjà  donné la réponse~: on
l'appelle $\ln 32$.

\end{comment}
%Voici un extrait des tables de logarithmes publiées au début du XVII%\eme
%siècle par John \textsc{Napier} himself :

%\begin{center}
%\vspace{-1cm}
%\includegraphics[width=0.7\linewidth]{napier_logtable}
%\end{center}


% 

%
%Plus généralement, nous allons voir si le mécanisme qui à $e^y$
%associe $y$ est une fonction et si oui découvrir ses propriétés
%qui seront bien sûr celle du logarithme de John.
%
%
%

\vspace{0.5cm}
\section{Fonction réciproque et théorème de la bijection}
Survolons donc la notion de \textit{fonction réciproque}. La
fonction exponentielle est continue et strictement croissante de
$\mathbb{R}$ vers $]0,+\infty[$. Donc, d'après notre bon vieux
\textsl{théorème de LA valeur intermédiaire} (appelé en d'autres
temps théorème de la bijection),  pour tout $x\in]0,+\infty[$, il
existe un unique réel $y$ tel que $\exp(y)=x$. Compte tenu de
l'existence et de l'unicité de ce réel, on peut donc
\textit{définir} une fonction qui à tout $x$ dans $]0,+\infty[$
associe l'unique réel $y$ tel que $\exp(y)=x$.\\
 
 C'est la fonction
réciproque de la fonction exponentielle, on la note $ln$ (logarithme
népérien)
\begin{center}ln~:
\begin{tabular}{l} $]0,+\infty[\to\mathbb{R}$\\
$ x\mapsto \ln x$
\end{tabular}
\end{center}

%Après ça, tout coule de source...







\section{Construction du logarithme népérien}


\begin{thm}[définition du logarithme népérien]
    La fonction exponentielle admet une
fonction réciproque défine sur $]0,+\infty[$ et à valeurs dans
$\bbr$ appelée fonction logarithme népérien. On note $\ln x$
l'image d'un réel strictement positif par cette fonction.
Alors,
\[\text{pour tout } x>0,\ e^{\ln x}=x \quad \text{et}\quad \text{pour tout réel }
x,\ \ln e^x=x\] 
    \end{thm}
    

\section{Conséquence sur la représentation graphique}
\begin{center}
\includegraphics{../Figures/expo.1}    
\end{center}
    
\section{Variation du logarithme népérien}

\begin{prop}[sens de variation]
   La fonction ln est dérivable sur
$]0,+\infty[$ et $\ln'(x)=\dfrac{1}{x}$ pour tout $x\in]0,+\infty[$. \\ On
en déduit que la fonction ln est strictement croissante sur
$]0,+\infty[$.
    \end{prop}

\begin{demo}
Formellement, on a $f \circ f^{-1} (x) = x$, dérivons cette relation, 
on obtient alors $(f^{-1})' (x) f'\circ f^{-1} (x)  =1 $ autrement dit, 
$(f^{-1})'(x) =\dfrac{1}{f' \circ f^{-1} (x)}$.

Posons $f(x)=e^x$ alors $f^{-1} (x) = ln x$ et $f'(x)=e^x$, par suite 
$f'\circ f^{-1} (x)=e^{lnx}=x$ d'où $ln x ' =\dfrac{1}{x}$.

Ainsi, si la fonction $ln$ est dérivable elle vaut nécéssairement $\dfrac{1}{x}$.

\end{demo}


\section{Propriétés du logarithme népérien}

\begin{prop}[relation fondamentale]
    Pour tout $a>0$ et tout $b>0$, on a
$$\ln ab = \ln a+\ln b$$
De plus $$\ln1=0$$
    \end{prop}




\begin{prop}[propriétés algébriques]
    $$\ln (1/a)=-\ln a \qquad
\quad\ln(a/b)=\ln a-\ln b\qquad \quad\ln a^p=p\ln a\quad
\text{avec } p\in\mathbb{Q}$$
    \end{prop}
\section{Limites du logarithme népérien et croissances comparées}


\begin{prop}[limite à l'infini]
     $$\limpinf  \ln x=+\infty$$
    \end{prop}


\begin{prop}[croissances comparées]
D'une part~:

    $$\limpinf \dfrac{\ln x}{x}=0$$

Plus généralement, pour tout entier naturel $n$ non nul $$\limpinf\fr{\ln x}{x^n}=0\quad \limpinf \fr{e^x}{x^n}=+\infty\quad
\limite{0}x^n\ln x=0\quad \limite{-\infty}x^ne^x=0$$

\end{prop}

\begin{demo}[preuve de $ \limite{0} x ln(x)$ ]
On sait que, $$ x ln(x) = \dfrac{1}{x} (- ln( \dfrac{1}{x} )) $$ 
ainsi, 
$$ \limite{0} x ln(x)= \limite{0} \dfrac{1}{x} (- ln( \dfrac{1}{x} ))$$
En posant $t=\dfrac{1}{x}$   on a quand $x$ tend vers 0 \emph{(à droite)} $t$ qui tend vers $+ \infty$, d'où, 
 $$ \limite{0} x ln(x)=\limite{0} \dfrac{1}{x} (- ln( \dfrac{1}{x} )) =- \limpinf \dfrac{ln(t)}{t} = 0^- $$ 
 
 
\end{demo}

\begin{prop}[limite en zéro et taux de variation]
    $$\limite{0}\frac{\ln(x+1)}{x}=1$$
    \end{prop}



\section{Logarithmes et exponentielles d'autres bases}


\begin{defn}[puissance réelle (exponentielles de base quelconque)]
    Pour tout réel $a>0$ et pour tout réel $b$, on pose $a^b=e^{b\ln a}$
    
    \vspace{3mm}
    \end{defn}
  

\begin{defn}[logarithme décimal]
    La fonction $\log\ :\ x\mapsto \fr{\ln x}{\ln 10}$ définie pour $x>0$ est appelée \textbf{fonction logarithme décimal}
    \end{defn}
  
  



\end{document}
